\begin{examplebox}
\textbf{\textcolor{CExample}{例一｜左右端点：公式比斜率更可靠}}

已知 $E:\dfrac{x^2}{9}+\dfrac{y^2}{4}=1$，求 $P(3,0)$ 处切线。

\textbf{验证：} $3^2/9+0^2/4=1$，确为椭圆上点。
直接代入：
\[
 \frac{3x}{9}+\frac{0\cdot y}{4}=1
 \quad\Longrightarrow\quad\boxed{x=3}.
\]
\textbf{点评：} 斜率不存在，但统一公式照样成立；绝不能在 $y_0=0$ 时使用斜率式。

\medskip
\textbf{\textcolor{CExample}{例二｜已知切点：求切线、截距和面积}}

已知 $E:\dfrac{x^2}{4}+\dfrac{y^2}{3}=1$，$P(1,\tfrac32)$，求切线及它与两坐标轴围成的三角形面积。

\textbf{第一步（验点）：}
\[
 \frac{1^2}{4}+\frac{(3/2)^2}{3}
 =\frac14+\frac34=1.
\]
\textbf{第二步（最短代入）：}
\[
 \frac{1\cdot x}{4}+\frac{(3/2)y}{3}=1
 \iff\boxed{\frac{x}{4}+\frac{y}{2}=1}.
\]
\textbf{第三步（截距和面积）：} 与坐标轴交于 $A(4,0)$ 和 $B(0,2)$：
\[
 S_{\triangle OAB}=\frac12\times 4\times 2=\boxed4.
\]
\textbf{点评：} 原例结论正确，本次补全了“验证切点”的关键步骤。

\medskip
\textbf{\textcolor{CExample}{例三｜反向题：已知斜率，直接求所有切线}}

椭圆 $E:\dfrac{x^2}{9}+\dfrac{y^2}{4}=1$，求斜率为 $1$ 的切线。

设 $y=x+m$。由 $m^2=a^2k^2+b^2$，
\[
 m^2=9\cdot1^2+4=13
 \quad\Longrightarrow\quad m=\pm\sqrt{13}.
\]
所以两条切线为
\[
 \boxed{y=x+\sqrt{13}},\qquad
 \boxed{y=x-\sqrt{13}}.
\]
对应切点也可秒求：
\[
 P=\left(-\frac{9}{m},\frac4m\right).
\]
\textbf{点评：} 已知斜率时，不必先猜切点；先求 $m$ 更快。

\medskip
\textbf{\textcolor{CExample}{例四｜一般式：怎样迅速判断是否为切线}}

判断 $2x+3y=t$ 与 $\dfrac{x^2}{9}+\dfrac{y^2}{4}=1$ 何时相切，并求 $t>0$ 的切点。

对于一般直线 $Ax+By=C$，若 $B\ne0$，化成斜截式：
\[
 k=-\frac AB,\qquad m=\frac CB.
\]
因此切线判定等价于（对于任意非零 $(A,B)$ 都成立）
\[
 \boxed{C^2=a^2A^2+b^2B^2}.
\]
代入 $A=2,B=3$：
\[
 t^2=9\cdot2^2+4\cdot3^2=72,
 \quad\boxed{t=\pm6\sqrt2}.
\]
当 $t=6\sqrt2$ 时，切点为
\[
 P=\left(\frac{a^2A}{C},\frac{b^2B}{C}\right)
 =\left(\frac{18}{6\sqrt2},\frac{12}{6\sqrt2}\right)
 =\boxed{\left(\frac3{\sqrt2},\sqrt2\right)}.
\]
验证：$(3/\sqrt2)^2/9+(\sqrt2)^2/4=1$，且 $2x_0+3y_0=6\sqrt2$。

\medskip
\textbf{\textcolor{CExample}{例五｜椭圆外一点：不能把外点当切点}}

从 $Q(4,0)$ 向椭圆 $E:\dfrac{x^2}{9}+\dfrac{y^2}{4}=1$ 引切线，求切线方程。

\textbf{第一步（判断）：} $4^2/9>1$，$Q$ 在椭圆外，不是切点。

\textbf{第二步（设过 $Q$ 的直线）：}
\[
 y=k(x-4)=kx-4k,
 \quad m=-4k.
\]
\textbf{第三步（判定相切）：}
\[
 m^2=9k^2+4
 \Longrightarrow16k^2=9k^2+4
 \Longrightarrow k=\pm\frac2{\sqrt7}.
\]
所以
\[
 \boxed{y=\frac2{\sqrt7}(x-4)},\qquad
 \boxed{y=-\frac2{\sqrt7}(x-4)}.
\]
\textbf{可选验算：} 两个切点分别为
$\left(\dfrac94,-\dfrac{\sqrt7}{2}\right)$ 与
$\left(\dfrac94,\dfrac{\sqrt7}{2}\right)$。
\textbf{点评：} $Q$ 的坐标直接代入 $x_0x/a^2+y_0y/b^2=1$ 绝不会得到“在 $Q$ 处的切线”，因为 $Q$ 根本不在椭圆上。

\medskip
\textbf{\textcolor{CExample}{高价值变式｜水平切线与竖直切线}}

同一椭圆 $x^2/9+y^2/4=1$：
\[
 k=0\Longrightarrow m^2=4\Longrightarrow y=\pm2;
 \qquad\text{竖直切线： }x=\pm3.
\]
\textbf{迁移要点：} $y=kx+m$ 只覆盖非竖直直线；“四个顶点、四条轴向切线”应能直接画出。
\end{examplebox}
